\subsection{带条带状端的纤维化}\label{sec:17b}
令 $M$ 为带角正合辛流形，$\widehat\pi:\widehat E\to\widehat S$ 为 Lefschetz 纤维化。假设底空间上有有限边界标点集 $\Sigma=\Sigma^-\sqcup\Sigma^+$，并在每个标点处将纤维 $\widehat E_\zeta$ 与 $M$ 认同。还假设纤维化在标点附近局部平凡。这意味着，每个 $\zeta$ 有邻域 $\widehat U=\widehat U_\zeta\subset\widehat S$，使
\[
\begin{aligned}
(&\widehat\pi|\widehat U:\widehat E|\widehat U\to\widehat U,
\omega_{\widehat E},\theta_{\widehat E},I_{\widehat E})\\
&=(\widehat U\times M\to\widehat U,
\omega_{\widehat S}+\omega_M,\theta_{\widehat S}+\theta_M,I_{\widehat S}\times I_M).
\end{aligned}
\]
照常取 $S=\widehat S\setminus\Sigma$；对每个 $\zeta$ 选择条带状端 $\epsilon_\zeta$，其像位于 $U_\zeta=\widehat U_\zeta\setminus\{\zeta\}$ 内。于是 $\widehat\pi$ 的限制 $\pi:E\to S$ 在各条带状端上带有典范平凡化 $Z^\pm\times M\to E$。称这样的对象为\emph{带条带状端的 Lefschetz 纤维化}。

\StatementNumber{6}\begin{note}\label{st:17.6}
许多情况下，标点附近的局部平凡性并不预先成立，但可如下通过变形补救。令 $\widehat\pi:\widehat E\to\widehat S$
\SourcePage{248}
为在 $\partial^h\widehat E$ 附近平凡的 Lefschetz 纤维化。选择 $\zeta\in\partial\widehat S$，置 $M=\widehat E_\zeta$。通过改变该纤维附近的辛结构及殆复结构，可在 $\zeta$ 附近达到局部平凡性；更准确地，用\secref{15b}的术语，这是沿纤维恒定的变形。与\stref{引理}{15.3}类似，在局部平凡化 $\widehat U\times M\to\widehat E|\widehat U$ 中有
\begin{equation}\label{eq:17.9}
\theta_{\widehat E}=\theta_M+\theta_{\widehat S}+\kappa+dR,
\end{equation}
其中 $\kappa$ 与函数 $R$ 均在 $\widehat U\times\partial M$ 附近消失。希望用 $\psi(\kappa+dR)$ 替换\eref{17.9}中的这两项，$\psi$ 是 $\widehat U$ 上在 $\zeta$ 附近为零的适当截断函数。修改后的一形式的外导数在总空间上未必仍为辛，但加上底空间某个适当项的拉回，就容易补救。类似地，在局部平凡化中，殆复结构 $I_{\widehat E}$ 在 $\widehat U\times\partial M$ 附近等于 $I_{\widehat S}\times I_M$；容易修改它，使其在整个纤维 $\widehat E_\zeta$ 附近都具有相同的乘积形式。
\end{note}

须将 Lagrangian 边界条件与相对扰动数据的概念扩展到条带状端背景。对前者，取 $F\subset E$，使其到底空间 $S$ 的投影为固有映射，并满足\eref{17.1}。则族 $F_z$ 在每个条带状端上自动局部恒定，即存在一对 Lagrangian 子流形 $(L_{\zeta,0},L_{\zeta,1})$，使对所有 $s$ 有 $F_{\epsilon_\zeta(s,k)}=L_{\zeta,k}$。此外，在条带状端上必有 $\alpha_F=0$。现假设已为每个 $(L_{\zeta,0},L_{\zeta,1})$ 固定 Floer 数据 $(H_\zeta,J_\zeta)$。则相对扰动数据仍为如前的一对 $(K,J)$，但另要求在每个条带状端上有 $K(s,t,x)=H_\zeta(t,x)\,dt$、$J(s,t,x)=i\times J_\zeta(t,x)$。取\eref{17.2}的解 $u$，利用每个端上的平凡化，写成 $(u\circ\epsilon_\zeta)(s,t)=(s,t,u_\zeta(s,t))$，其中 $u_\zeta:Z^\pm\to M$。于是\eref{17.2}归结为 $u_\zeta$ 的 Floer 方程，故可合理要求其渐近收敛于某个 $y_\zeta\in\mathcal C(L_{\zeta,0},L_{\zeta,1})$。\znote{240}{条带状端背景下的三个技术要点值得串起来看。一、“投影固有”不可省：$F_z$ 须有界地覆盖到 $S$ 的端点，端上的极限 $L_{\zeta,k}$ 才存在。二、端上 $F$ 自动局部恒定，故每条端只记录一对 $(L_{\zeta,0},L_{\zeta,1})$，$\alpha_F$ 也随之在端上为零；这一步同样依赖投影固有性。三、条件 $K=H_\zeta(t,x)\,dt$、$J=i\times J_\zeta$ 是让\eref{17.2}在端上恰好退化为 Floer 方程的自然选择，于是每条解在端上有极限 $y_\zeta$，计数才落在 $CF^{\mathrm{pr}}$ 上。[来源：原书 §(17b)，印刷页 240–241。]}所得模空间记为 $\M_{E/S}(\{y_\zeta\})$。假设满足适当正则性条件，就可如\eref{8.20}，对 $\M_{E/S}(\{y_\zeta\})$ 中孤立点作模二计数，定义链同态
\begin{equation}\label{eq:17.10}
C\Phi_{E/S}:\bigotimes_{\zeta^+\in\Sigma^+}CF^{\mathrm{pr}}(L_{\zeta^+,0},L_{\zeta^+,1})
\longrightarrow\bigotimes_{\zeta^-\in\Sigma^-}CF^{\mathrm{pr}}(L_{\zeta^-,0},L_{\zeta^-,1}).
\end{equation}
在 Floer 上同调上诱导的映射 $\Phi_{E/S}$ 与扰动数据的选择无关，并满足 TQFT 公理的直接推广。特别地，有一个复合律：将两个 Lefschetz 纤维化 $\pi_k:E_k\to S_k$，通过把 $S_1$ 的一个出向端与 $S_2$ 的一个入向端认同来粘合。当然，$S_k$ 上辛结构通常不一致，但可如\stref{笔记}{17.6}，通过局部拼接处理。

\SourcePage{249}
这一框架既包含前面考虑的 $S$ 紧致的情形，也包含第\sr{sec:8}{8 节}的理论；后者对应于只允许平凡纤维化 $E=S\times M$。借用这些情形中的想法，还可沿若干方向进一步扩展：
\begin{itemize}
\item \emph{移动边界条件。}如紧致情形中所讨论，但另须要求 $F_z$ 在每个条带状端上关于 $z$ 局部恒定。
\item \emph{带参数的模空间。}在基本构造中已经出现 Lefschetz 纤维化族及相应模空间，用它们为不同相对扰动数据选择所给出的映射\eref{17.10}构造链同伦。一般结构是：考察依赖于参数空间中 $r\in R$ 的族 $E^r\to S^r$，带 Lagrangian 边界条件 $F^r$ 和扰动数据 $(K^r,J^r)$，而条带状端结构关于 $r$ 恒定。为简单起见，假设 $R$ 紧致。相应的 Floer 上链群之间的映射，满足形如 $\partial\circ(\text{映射})+(\text{映射})\circ\partial+(\text{边界项})=0$ 的等式，其中边界项来自将族限制到 $\partial R$。特别地，$R=[0,1]$ 时，得到对应于两个端点的两映射之间的链同伦。后面会看到几个双参数例子；进一步说，可以沿第\sr{sec:9}{9 节}的思路建立一般理论，但对所得代数结构的意义，目前只理解了一部分。
\item \emph{求值映射。}除条带状端外，还可在 $\partial S$ 上固定有限点集 $\Pi=\Pi^-\sqcup\Pi^+$，并考虑这些点处的求值映射 $\operatorname{ev}_{z^\pm}(u)=u(z^\pm)$。由此得到的不变量，如\eref{8.24}，在张量积中具有额外因子 $H^*(F_{z^\pm};\K)$。例如，在\stref{例题}{17.3}的局部模型中，任意点处的求值映射可认同为从切球面丛到底空间的投影 $S(TS^n)\to S^n$。从根本上说，引入标点并不使理论更一般，因为处理技术细节并使用 PSS 同构后，可用条带状端替代这些标点；但它在计算中往往更方便。对此情形还有复合律的一种变体：沿标点取底曲面的边界连通和。
\end{itemize}

\subsection{一个消失定理}\label{sec:17c}
最后提到的复合律，会在后续讨论中多次使用，因此值得说得更准确些。为简单起见，令 $\pi_k:E_k\to S_k$，$k=1,2$，为两个 Lefschetz 纤维化，第一个带条带状端，第二个具有紧底空间，各自带有可能移动的 Lagrangian 边界条件 $F_k$。此外，还要有边界标点 $z_1\in\partial S_1$、$z_2\in\partial S_2$，使纤维 $(E_k)_{z_k}$
\SourcePage{250}
都可与共同的 $M$ 认同，且 $(F_k)_{z_k}$ 成为同一子流形 $L\subset M$。还假设两个纤维化在标点附近局部平凡，边界条件局部恒定。选择相对扰动数据 $(K_k,J_k)$，在 $(E_k)_{z_k}$ 附近平凡，即 $K_k=0$、$J_k=I_{E_k}$。现在在标点处作 $S_1,S_2$ 的边界连通和。可从两种方式理解：一种是先去掉 $z_k$，在所得穿孔曲面上引入条带状端，再对某个粘合长度 $l$，如\eref{8.5}进行端的标准粘合；另一种是在 $z_k$ 附近取以半圆盘 $D_+=D\cap H$ 为模型的局部坐标，再对某个 $\rho\in(-1,0)$ 使用形如 $z\mapsto\rho z^{-1}$ 的过渡映射。两种图景完全等价，具有通常的关系 $l=-\log(-\rho)/\pi$。\znote{242}{原书此处印作 $l=-\log(\rho)/\pi$，但同句已声明 $\rho\in(-1,0)$，负实数在实对数下无定义，故此式只能按 $-\log|\rho|/\pi$ 读，译文径写 $-\log(-\rho)/\pi$。因子 $\pi$ 而非 $2\pi$ 来自模型：标点附近取的是半圆盘 $D_+=D\cap H$，其角为 $\pi$，而 $z\mapsto\rho z^{-1}$ 是穿过该角的过渡映射；绕标点一周的“角度”与 $|\rho|$ 的对数成正比，比例即半圆的 $\pi$。相应地，$|\rho|$ 越小，粘合长度 $l$ 越大，这与 $\rho\to0$ 对应长粘合的说法一致。[来源：原书 §(17c)，印刷页 242；对数定义域为初等核对。]}提及它们，是因为两者各有直观性，并暗示两个略有不同但同样可行的分析框架：第一种得到通常 Floer 理论的 Morse–Bott 版本，第二种接近相对 Gromov–Witten 理论中全纯圆盘的粘合程序。不论哪一种，所得空间均带有 Lefschetz 纤维化 $\pi:E\to S$，以及诱导的边界条件 $F$ 和相对扰动数据 $(K,J)$；基本粘合定理\bcite{131}等式 (2.28) 说明，若满足适当正则性条件，且 $l$ 大，或 $\rho$ 接近 $0$，则
\begin{equation}\label{eq:17.11}
\M_{E/S}(\{y_\zeta\})^0\cong
\coprod_{p+q=n}\M_{E_1/S_1}(\{y_\zeta\})^p\times_L\M_{E_2/S_2}^q.
\end{equation}
这里上标表示维数：右边的维数如此安排，使关于求值映射 $\operatorname{ev}_{z_1}:\M_{E_1/S_1}(\{y_\zeta\})\to(F_1)_{z_1}\cong L$ 及 $\operatorname{ev}_{z_2}:\M_{E_2/S_2}\to(F_2)_{z_2}\cong L$ 所取的纤维积为零维。上述正则性假设包括 $\operatorname{ev}_{z_1},\operatorname{ev}_{z_2}$ 相互处于一般位置。现在将其转成相应不变量的结论。对 $E_2$，考虑由 $\operatorname{ev}_{z_2}$ 表示的同调类，记为 $\Phi_{E_2/S_2,z_2}\in H_*(L;\K)$。对 $E$，则直接取\eref{17.10}。

\StatementNumber{7}\begin{lemma}\label{st:17.7}
若 $\Phi_{E_2/S_2,z_2}$ 消失，则 $\Phi_{E/S}$ 也消失。
\end{lemma}
\begin{proof}
证明直接：取以 $\operatorname{ev}_{z_2}$ 为边界的（伪）循环 $b:B\to L$。则 $\M_{E_1/S_1}(\{y_\zeta\})\times_L B$ 的零维部分给出 Floer 上链群之间的映射，它是 $\Phi_{E/S}$ 的零同伦。
\end{proof}

\clearpage
\StatementNumber{8}\begin{note}\label{st:17.8}
我们对该引理的所有应用中，紧的部分都如\stref{引理}{17.4}。此时可加上假设 $2c_1(M,L)=0$，并得到稍强的结果。若 $n>1$，空间 $\M_{E_2/S_2}$ 的维数为 $2n-1>n$，所以纤维积\eref{17.11}为空。在余下的 $n=1$ 情形，唯一非平凡的项为 $p=0,q=1$。由\stref{引理}{17.4}，已知求值映射
\[
\operatorname{ev}_{z_2}:\M_{E_2/S_2}^1\longrightarrow L
\]
\SourcePage{251}
具有偶数次。对每个孤立点 $u_1\in\M_{E_1/S_1}(\{y_\zeta\})^0$，由于横截性假设，$\operatorname{ev}_{z_1}(u_1)$ 是 $\operatorname{ev}_{z_2}$ 的正则值，故满足 $\operatorname{ev}_{z_2}(u_2)=\operatorname{ev}_{z_1}(u_1)$ 的映射 $u_2$ 的数目为偶数。这说明纤维积仍包含偶数个点。结论是，对这一特定的扰动数据选择，上链层次的不变量 $C\Phi_{E/S}$ 本身就是零。
\end{note}
